\documentclass[11pt]{article}
\usepackage[margin=1in]{geometry}
\usepackage{amsmath,amssymb,amsthm,hyperref,enumitem}
\newtheorem{theorem}{Theorem}
\newtheorem{lemma}{Lemma}
\title{The sine-angle invariant of harmonic polygons}
\author{Principia Math In Progress}
\date{21 September 2026}
\begin{document}
\maketitle
\noindent\textbf{MathDB 352834.} This note proves Conjecture 2 of
Garcia, Reznik, and Roitman \cite[p.~228]{GRR}, for every number of sides.
It concerns internal angles of a single harmonic polygon and is distinct
from their Conjecture 1 about two lateral polar areas. Review is by the
authoring agent only.

\section*{Task 1: the full sine-angle invariant}
Write $A$ for the area, $s_j$ for the sidelengths, and $\theta_j$ for the
internal angles of a convex harmonic $N$-gon, $N\ge3$. The standard
parametrization of its family, up to a fixed Euclidean similarity, is
\begin{equation}\label{vertices}
 P_j=O+Re^{i\gamma}w_j,\qquad
 w_j=\frac{z_j-d}{1-dz_j},\qquad
 z_j=e^{i(t+j\beta)},\qquad \beta=\frac{2\pi}{N},
\end{equation}
where $R>0$, $-1<d<1$, and $O,\gamma$ are fixed; only $t$ varies.
This is the construction in \cite[Section 2]{GRR}, with an immaterial
rotation by $\pi$ absorbed into $\gamma$. It includes all the harmonic
families in the source. Let
\[
 D_m=1-2d^2\cos(m\beta)+d^4>0,\qquad
 S=\sum_{j=0}^{N-1}\sin(2\theta_j),\qquad Q=\sum_{j=0}^{N-1}s_j^2.
\]

\begin{theorem}\label{main}
For every $N\ge3$ and every parameter in \eqref{vertices},
\begin{align}
 \frac{S}{A}&=-\frac{2\sin(2\beta)D_1}{R^2\sin\beta\,D_2},\label{ratioA}\\
 \frac{S}{Q}&=-\frac{(1-d^2)\sin(2\beta)D_1}
 {2R^2(1+d^2)(1-\cos\beta)D_2}.\label{ratioQ}
\end{align}
Both ratios are independent of $t$. These formulas include $d=0$ and
$N=4$, when $S=0$.
\end{theorem}

The key point is that the correlations for different vertex steps have
the same phase factor. Put
\[
 H(t)=\frac{1-d^{2N}}{1-2d^N\cos(Nt)+d^{2N}}>0,\qquad
 T_m=\sum_{j=0}^{N-1}\overline{w_j}w_{j+m},
\]
with indices taken modulo $N$.
\begin{lemma}\label{trace}
For every integer $m$,
\begin{equation}\label{traceformula}
 T_m=N(1+C_mH(t)),\qquad
 C_m=\frac{(1-d^2)(e^{im\beta}-1)}{1-d^2e^{im\beta}}.
\end{equation}
In particular,
\begin{equation}\label{parts}
 \operatorname{Im}C_m=\frac{(1-d^2)^2\sin(m\beta)}{D_m},\qquad
 \operatorname{Re}C_m=\frac{(1-d^4)(\cos(m\beta)-1)}{D_m}.
\end{equation}
\end{lemma}
\begin{proof}
Set $v=e^{im\beta}$. Since $|w_j|=1$, the summand equals
\[
 \frac{(vz_j-d)(1-dz_j)}{(1-dvz_j)(z_j-d)}
 =1+C_m\left(\frac1{1-dvz_j}+\frac{d}{z_j-d}\right).
\]
This is a rational identity, verified by multiplying the denominators.
Multiplication by $v$ permutes the $z_j$. The root-of-unity identity
$\sum_jz_j^k=0$ for $N\nmid k$, applied to convergent geometric series,
therefore gives
\[
 \sum_j\frac1{1-dvz_j}=\frac{N}{1-d^Ne^{iNt}},\qquad
 \sum_j\frac{d}{z_j-d}=\frac{Nd^Ne^{-iNt}}{1-d^Ne^{-iNt}}.
\]
Their sum is $NH(t)$, proving \eqref{traceformula}. Multiplying the
numerator and denominator of $C_m$ by $1-d^2e^{-im\beta}$ gives
\eqref{parts}. All denominators are nonzero for $|d|<1$.
\end{proof}

\begin{proof}[Proof of Theorem \ref{main}]
The map in \eqref{vertices} preserves cyclic order on the circle.
Choose increasing lifts of the arguments $\phi_j$ of $w_j$, with
$\phi_N=\phi_0+2\pi$, and put $\delta_j=\phi_{j+1}-\phi_j>0$.
For the convex cyclic polygon, the inscribed-angle theorem gives
\[
 \theta_j=\pi-\frac{\delta_{j-1}+\delta_j}{2}.
\]
Consequently,
\begin{equation}\label{angletrace}
 S=-\sum_j\sin(\delta_{j-1}+\delta_j)
   =-\operatorname{Im}\sum_j\overline{w_{j-1}}w_{j+1}
   =-\operatorname{Im}T_2.
\end{equation}
The ordinary positive area, by the shoelace formula, is
\begin{equation}\label{areatrace}
 A=\frac{R^2}{2}\operatorname{Im}T_1.
\end{equation}
Also $|w_j|=1$ implies
\begin{equation}\label{sidetrace}
 Q=R^2\sum_j|w_{j+1}-w_j|^2
  =2R^2(N-\operatorname{Re}T_1).
\end{equation}
Lemma \ref{trace} now yields the three exact formulas
\[
 S=-NH(t)\operatorname{Im}C_2,\qquad
 A=\frac{NR^2}{2}H(t)\operatorname{Im}C_1,\qquad
 Q=-2NR^2H(t)\operatorname{Re}C_1.
\]
Here $H(t)>0$, $\operatorname{Im}C_1>0$, and
$\operatorname{Re}C_1<0$, since $0<\beta<\pi$.
Division is therefore legitimate. Cancelling the common factor and
substituting \eqref{parts} proves both \eqref{ratioA} and \eqref{ratioQ}.
If $N=4$, $\sin(2\beta)=0$ gives the zero numerator without any division
by it. If $d=0$, all formulas remain valid directly.
\end{proof}

\noindent\textbf{Attribution and distinction.}
The family parametrization and known area--squared-sidelength invariant
are prior results in \cite{GRR}. Both requested ratios were nevertheless
derived directly here. The root-of-unity correlation method extends the
area calculation in our immediately preceding note for MathDB 352833.
That note proves the distinct two-focus reciprocal-area assertion; it
does not state or prove the internal-angle assertion. This note adds the
step-two trace and its geometric angle interpretation. It earns one
credit for Conjecture 2, not separate credits for the two equivalent ratios.

\section*{Tasks 2--5: proposed follow-ups}
The following are proposals, not additional completed solves.
\begin{enumerate}[label=\textbf{Task \arabic*.},start=2,leftmargin=*]
\item Classify the ratios of sums of trigonometric polynomials in internal
angles to area that are invariant for every harmonic family.
\item Express the constants in \eqref{ratioA} and \eqref{ratioQ} directly
in terms of the Brocard angle and circumradius, eliminating $d$.
\item Determine which products of nonadjacent edge directions have
correlation sums with a common phase factor.
\item Extend the angle identities to harmonic star polygons, with explicit
conventions for directed angles and signed area.
\end{enumerate}

\begin{thebibliography}{9}
\bibitem{GRR} R. Alves Garcia, D. Reznik, and P. Roitman,
\emph{New Properties of Harmonic Polygons}, Journal for Geometry and
Graphics \textbf{26}(2) (2022), 217--236.
\href{https://www.heldermann-verlag.de/jgg/jgg26/j26h2alve.pdf}{Publisher PDF}.
Current preprint: \href{https://arxiv.org/abs/2112.02545}{arXiv:2112.02545v3}.
Both complete versions were read; Conjecture 2 remains stated on p.~228.
\end{thebibliography}
\end{document}
